\chapter{Rangkuman}

\section{Rangkuman Pelaksanaan Pekerjaan}
Kerja praktik dilaksanakan di PT Radya Anugrah Digital pada 22 Juni sampai dengan 22 September 2026 dengan penempatan sebagai \textit{Quality Assurance} pada proyek Komatsu ESS. Pekerjaan berjalan dalam enam tahap, yaitu orientasi, penyusunan draf skrip uji, pengujian per modul, pengujian antarmodul, pengujian regresi, dan dokumentasi. Pada setiap tahap pengujian, cacat yang ditemukan dilaporkan kepada tim pengembang dan perbaikannya diverifikasi. Pekerjaan dilakukan bersama seorang rekan KP dengan pembagian modul yang ditetapkan pembimbing, dan kegiatan harian dicatat pada \textit{log activity}.

\section{Rangkuman Hasil Pekerjaan}
Hasil pekerjaan penulis berupa skrip uji, laporan cacat, dan laporan hasil pengujian yang digunakan tim pengembang untuk meningkatkan kualitas aplikasi Komatsu ESS. \TEKNIS{ringkasan hasil kuantitatif dan pekerjaan yang belum selesai} Melalui pekerjaan ini penulis memperoleh kompetensi berikut:
\begin{enumerate}[label=\arabic*)]
\item menyusun skrip uji dari kebutuhan fitur;
\item menguji aplikasi secara bertahap dari modul hingga lintas modul dan regresi;
\item melaporkan cacat secara jelas dan dapat direproduksi; dan
\item bekerja dalam tim, berkomunikasi dengan pengembang, dan menjaga kerahasiaan informasi.
\end{enumerate}

\section{Evaluasi Diri}
Penulis menilai bahwa kerja praktik ini memberi gambaran nyata tentang peran penguji dalam pengembangan perangkat lunak yang tidak diperoleh dari perkuliahan. Penulis berhasil menyelesaikan tahapan pengujian sesuai rencana dan menerapkan sikap profesional, seperti tepat waktu dan menjaga kerahasiaan. Di sisi lain, pada awal KP penulis membutuhkan waktu untuk memahami alur bisnis aplikasi sehingga penyusunan skrip uji berjalan lebih lambat. Ke depan, penulis perlu lebih awal bertanya dan mendokumentasikan pemahaman kebutuhan sebelum menyusun skrip uji. \ISI{sesuaikan evaluasi diri dengan pengalaman pribadi: keberhasilan, kegagalan, dan hal yang ingin diperbaiki}

\section{Ucapan Terima Kasih}
Penulis mengucapkan terima kasih kepada Bapak Hari Purnama selaku dosen pembimbing, Ibu Alyssa Christella selaku \textit{Project Manager} sekaligus \textit{mentor} di PT Radya Anugrah Digital, Bapak Hari Bagus Firdaus selaku Direktur, seluruh tim proyek Komatsu ESS, serta Zulfaqqar Nayaka Athadiansyah selaku rekan KP, atas bimbingan dan bantuan selama pelaksanaan kerja praktik.
